\headright{Професионален опит}

\jobtitle{DevOps инженер}{Alten (Франция)}{05.2025 -- настояще}{}
\small

\noindent \textbf{Amadeus}

\begin{itemize}[label=\tiny$\bullet$, leftmargin=1em]
	\item Мигрирах стари Groovy скриптове за пипалини към вътрешна рамка, която управлява пълния жизнен цикъл на приложенията.
	\item Пренаписах рамката за фийдове от \textbf{Bash} на \textbf{Python} за планирани доставки на данни.
	\item Въведох индустриални рамки за автоматизиране на локалните зависимости (Localstack, Nix devenv).
	\item Разработих скриптове за динамично превключване на типовете машини в EMR клъстерите от спот към по заявка въз основа на латентността на опашката.
	\item Преработих вътрешната документация за по-голяма яснота и лесно поддържане.
	\item Проектирах HLD за управление на версиите, стандартизирайки работните потоци за внедряване в съответствие с изискванията за управление на промените.
\end{itemize}

\vspace{3.0ex}
\jobtitle{DevOps инженер}{DXC Technology (България)}{02.2022 -- 04.2025} {}
\small

\noindent \textbf{Continental AG (Гуми)}
\begin{itemize}[label=\tiny$\bullet$, leftmargin=1em, nosep]
	\item Ръководих екип от 5 души и докладвах директно за състоянието на проекта на клиента.
	\item Оптимизирах логиката за обработка на данни чрез мигриране на стари Spark задачи за пакетна обработка към Apache Flink стрийминг пипалини, което значително намали латентността.
	\item Ускорих стартирането на проекта чрез разрешаване на критични проблеми с VPN свързаността и автоматизиране на локалните процедури за разработка.
\end{itemize}

\vspace{1em}
\noindent \textbf{Nestle}
\begin{itemize}[label=\tiny$\bullet$, leftmargin=1em, nosep]
	\item Внедрих автоматизирани проверки за структурна валидация на Pull Requests, за да предотвратя несъответстващи разгръщания.
	\item Разширени диалекти на Snowflake в SQLFluff и автоматизирана интеграция в етапа преди пускане
	\item Интегрирани ANTLR парсери в \textbf{CI/CD} пипалината като задължителна проверка за валидиране преди пускане.
\end{itemize}

\vspace{1em}
\noindent \textbf{Ahold Delhaize}
\begin{itemize}[label=\tiny$\bullet$, leftmargin=1em]
	\item Препроектирах рамката за автоматизация на тестовете в модулна архитектура, подобрявайки възможността за повторна употреба на кода и намалявайки времето за въвеждане на QA инженерите.
	\item Сътрудничих с QA екипа в седмични синхронизации за идентифициране на оперативни затруднения, което доведе до подобрена автономност на екипа.
	\item Проектирах HLD за споделяне на данни между различни среди, което подкрепи миграцията на новата рамка за въвеждане в облака (CF2).
	\item Автоматизирах поток за хостинг на Allure отчети като статичен уебсайт в Azure Blob Storage с персонализиран мидълуеър за удостоверяване.
	\item Ръководих техническа оценка на \textbf{SonarQube} за автоматизиран анализ на код, с цел да се идентифицират изискванията за интеграция на платформата
\end{itemize}

\vspace{1em}
\noindent \textbf{Continental AG (ADAS, кодово име CAEdge)}
\begin{itemize}[label=\tiny$\bullet$, leftmargin=1em]
	\item Автоматизирах стандартизирана \textbf{VPC} топология за AWS конфигурация с множество акаунти, като я интегрирах в оркестратора Account Vending Machine.
	\item Проектирах и автоматизирах решение за управление на разходите за облака в организация с множество AWS акаунти, като внедрих бюджетни предупреждения за предотвратяване на превишения.
	\item Намалих повърхността за атака на облачните приложения чрез отстраняване на CVE, идентифицирани от докладите на AWS Inspector.
\end{itemize}

\vspace{5.0ex}
\jobtitle{Стажант софтуерен инженер}{Sys Consulting Ltd.}{07.2021 -- 12.2021}{}
\small
Поддържах стари банкови системи в CAVO за Windows XP среди.
